\begin{abstract}

Traditional HVAC systems, designed with a "one-size-fits-all" approach, often fail to provide thermal comfort in modern office environments, which are characterized by increasingly flexible and dynamic usage.
Occupant-centric control (OCC) has emerged as a promising solution to enhance both occupant satisfaction and energy efficiency by personalizing thermal conditions.
However, its widespread real-world adoption is hindered by two critical challenges: (1) the inability of existing frameworks to effectively manage the highly dynamic and unpredictable occupancy patterns typical of post-pandemic workplaces, and (2) the prohibitive survey effort required to build and maintain accurate personal comfort models (PCMs) for every individual.

This research aims to address these two gaps by developing and validating a practical OCC framework designed for real-world, multi-occupant office environments.
The methodology consists of a two-part field experiment conducted in an operational office building.
The first experiment will systematically investigate the relationship between various occupancy patterns—such as occupancy variability and user turnover—and the performance of dynamic OCC.
Its objective is to identify the specific conditions under which real-time, personalized control offers significant advantages over conventional static control strategies.

Building on these findings, the second experiment will propose and evaluate a novel, low-effort framework for the operational development of comfort models.
This framework employs an active transfer learning approach to intelligently minimize the total survey burden by strategically deciding when, where, and from whom to collect comfort feedback.
The system prioritizes queries that are most informative for improving group comfort and unlocking energy-saving potential, thereby balancing data efficiency with operational effectiveness.

The anticipated outcome of this research is a set of data-driven, actionable guidelines for building managers on selecting and deploying appropriate OCC strategies based on workspace usage patterns.
Furthermore, the proposed low-effort framework will provide a scalable and economically viable solution for implementing advanced, personalized control in commercial buildings.
By bridging the gap between academic theory and practical application, this work will contribute significantly to creating comfort and energy-efficient indoor environments tailored to the demands of modern work styles.

\end{abstract}
